% ========== 前言 ==========
\chapter*{前言}
\addcontentsline{toc}{chapter}{前言}

\subsection*{研究背景与动机}

中国正经历快速人口老龄化，60岁及以上人口已超2.8亿。独居老人的安全监护成为社会迫切需求——跌倒是65岁以上老人意外伤害的首要原因，每年约30\%的老人至少跌倒一次。然而，单一模态的监护方式各有局限：视频监控存在隐私争议且在遮挡弱光下失效，音频监测误报率高，穿戴式生理设备老人依从性差，用药记录单独无法反映即时风险。

多模态融合能够互补各模态优势，提升风险评估的可靠性与鲁棒性。本工作提出``智护家''系统，基于跨模态注意力机制（Cross-Modal Attention）融合视频、音频、生理数据、用药信息四种模态，进行三级风险评估（低/中/高风险），并开发完整的Web应用，支持缺失数据处理、模态贡献可视化和健康周报生成。

\subsection*{核心问题}

本工作聚焦四个核心问题：（1）如何将异构的四模态数据有效融合，而非简单拼接？（2）跨模态注意力机制能否学习到模态间的交互关系？（3）在真实公开数据上训练的融合模型，各模态贡献如何量化？（4）模态数据缺失时系统如何保持可用性？

\subsection*{主要成果}

系统在真实公开数据集上训练：ESC-50环境声音（2000条音频）提供真实声学特征，Pexels公开视频（112段）经MediaPipe PoseLandmarker提取人体骨架运动特征。Full模型（8.5M参数）在90样本独立测试集上达到100\%准确率，消融实验验证四模态全部参与决策（健康数据最重要，下降32.22\%；视频/音频各贡献1.11\%）。端到端真实音频推理6/6正确。

\subsection*{报告组织}

第\ref{chap:intro}章介绍研究背景与贡献；第\ref{chap:related}章回顾相关工作；第\ref{chap:method}章描述系统架构与方法；第\ref{chap:data}章介绍数据集与特征提取；第\ref{chap:exp}章给出实验设置与结果分析；第\ref{chap:discuss}章讨论局限性与未来工作；第\ref{chap:concl}章总结全文。
